\item Un rayo de luz entra a la atmósfera de la Tierra y desciende verticalmente a la superficie que está a una distancia $h = 100\text{ km}$ hacia abajo. El índice de refracción donde la luz entra a la atmósfera es de $1.00$ y aumenta linealmente con la distancia a la superficie del planeta, donde tiene un valor de $n = 1.000293$ en la superficie de la Tierra.
\begin{enumerate}
    \item ¿Cuánto tarda el rayo en recorrer esta trayectoria?
    \item ¿En qué porcentaje es mayor este intervalo de tiempo que el necesario en la ausencia de la atmósfera de la Tierra?
\end{enumerate}